\section{问题二模型的建立与求解}
【此处插入图：问题二全年逐日购电结构】
【此处插入图：问题二日末储电量轨迹】
【此处插入图：问题二指定日期计划购电与紧急购电】

\subsection{数据处理与建模思路}

问题二与问题一相比，最大的变化并不是微网物理结构发生变化，而是每天 0:00 制定计划时无法提前知道当天真实负荷和真实光伏。若全天计划不足，则需要按照交易时刻电价的 5 倍进行紧急购电。因此，问题二必须区分“0:00 事前计划”和“实际运行后的物理执行”。

问题一中已经建立了 \(G^L,G^{ch},PV^{ch},C,D,E,V,u\) 组成的能量分流物理模型。问题二继续沿用这组变量，但这些物理执行变量不应在 0:00 全部锁死。题目真正要求提前确定的是全天计划购电量，因此增加第一阶段变量

\[
G^{plan}_{d,t},
\]

表示日期 \(d\) 在 0:00 对第 \(t\) 个 10 min 时段提交的正常计划购电功率。

实际负荷和光伏实现后，如果计划电和储能仍无法满足负荷，则需要紧急购电，因此增加

\[
H_{d,t},
\]

表示紧急购电功率。

另一方面，计划购电已经在 0:00 确定并按计划量计费，当实际负荷低于预测、储能又无法继续吸收时，计划电中可能有一部分无法利用，因此增加

\[
W_{d,t},
\]

表示已购未用电功率。由此，问题二的变量并不是重新设计，而是在问题一物理变量之上由题目的计划机制自然增加 \(G^{plan},H,W\)。

每天仍划分为

\[
T=144,\qquad \Delta t=\frac16\ {\rm h}.
\]

2025 年 1 月作为已知数据预热期，直接使用真实负荷和真实光伏建立预测历史；正式滚动规划从 2025 年 2 月 1 日开始。

储能参数保持

\[
1200\le E_{d,t}\le10800,
\]

\[
0\le C_{d,t},D_{d,t}\le5000,
\]

\[
\eta_c=\eta_d=0.90.
\]

问题二不再要求每天首末储电量相同，而采用跨日连续：

\begin{equation}
E^{act}_{d+1,0}=E^{act}_{d,T}.
\end{equation}

\subsection{完美信息确定性基准模型}

为了衡量预测不确定性本身带来的额外成本，首先建立逐日完美信息基准。假设每天 0:00 已经知道当天真实负荷和真实光伏，则可在问题一模型基础上增加紧急购电变量 \(H\)。

定义

\[
PV^L_{d,t}=\min(PV_{d,t},L_{d,t}),
\]

\[
\overline L_{d,t}
=
\max(L_{d,t}-PV_{d,t},0),
\]

\[
\overline{PV}_{d,t}
=
\max(PV_{d,t}-L_{d,t},0).
\]

负荷平衡为

\begin{equation}
G^L_{d,t}+D_{d,t}+H_{d,t}
=
\overline L_{d,t}.
\end{equation}

储能充电来源为

\begin{equation}
PV^{ch}_{d,t}+G^{ch}_{d,t}
=
C_{d,t},
\end{equation}

光伏剩余平衡为

\begin{equation}
PV^{ch}_{d,t}+V_{d,t}
=
\overline{PV}_{d,t}.
\end{equation}

储能状态转移为

\begin{equation}
E_{d,t}
=
E_{d,t-1}
+
0.9C_{d,t}\Delta t
-
\frac{D_{d,t}\Delta t}{0.9}.
\end{equation}

充放电互斥继续通过

\[
C_{d,t}\le5000u_{d,t},
\]

\[
D_{d,t}\le5000(1-u_{d,t}),
\]

\begin{equation}
u_{d,t}\in\{0,1\}
\end{equation}

实现。

逐日完美信息基准的目标函数为

\begin{equation}
\min Z^{PI}
=
\sum_d\sum_t
\left[
\pi_t(G^L_{d,t}+G^{ch}_{d,t})
+
5\pi_tH_{d,t}
\right]\Delta t.
\end{equation}

由于此时当天真实负荷和光伏完全已知，且正常购电量没有题面给定上限，因此理论上可以在正常价格下提前购买足够电量，故基准模型应得到

\[
H_{d,t}=0.
\]

该模型不是现实可执行方案，而是用于衡量“如果拥有完美信息，成本最低可以达到多少”。

\subsection{7 日滚动 SAA 两阶段随机 MILP}

现实情况下，每天 0:00 只能根据历史数据预测未来负荷和光伏，因此正式模型采用两阶段随机规划。

\subsubsection{第一阶段：日前计划变量}

每天 \(d\) 的 0:00 确定

\begin{equation}
G^{plan}_{d,t}\ge0,
\qquad t=1,\ldots,144.
\end{equation}

该变量在所有 SAA 情景中完全相同，体现“计划在随机情景实现前已经提交”的非预见性。

\subsubsection{第二阶段：情景实现后的执行变量}

对每个随机情景 \(\omega\)，定义

\[
G^{L,(\omega)},\quad
G^{ch,(\omega)},\quad
PV^{ch,(\omega)},\quad
C^{(\omega)},\quad
D^{(\omega)},\quad
E^{(\omega)},\quad
V^{(\omega)},\quad
H^{(\omega)},\quad
W^{(\omega)},\quad
u^{(\omega)}.
\]

其中 \(H\) 表示紧急购电，\(W\) 表示计划购电中已经支付但未被负荷或储能利用的部分。

计划与执行之间满足

\begin{equation}
G_{d,t}^{L,(\omega)}
+
G_{d,t}^{ch,(\omega)}
+
W_{d,t}^{(\omega)}
=
G_{d,t}^{plan}.
\end{equation}

这样可以保证计划购电量不会在实际运行时无成本消失。

\subsection{预测模型与 SAA 情景生成}

每天 \(d\) 只使用 \(d\) 日以前已经发生的历史真实数据建立负荷和光伏中心预测，记为

\[
\hat L_{\tau,t}^{(d)},
\qquad
\widehat{PV}_{\tau,t}^{(d)}.
\]

在基础点预测上保留偏差校正，以降低系统性供电不足风险。

对历史日期 \(r\)，定义负荷和光伏预测误差

\[
e^L_{r,t}
=
L_{r,t}^{act}
-
\hat L_{r,t},
\]

\begin{equation}
e^{PV}_{r,t}
=
PV_{r,t}^{act}
-
\widehat{PV}_{r,t}.
\end{equation}

每个历史日保留完整的 144 槽误差向量，负荷误差与光伏误差尽量从同一历史日期联合抽取，不逐槽独立打乱。

正式模型保留最近 28 个有效历史残差日，并随机无放回抽取

\[
K=4
\]

个完整误差日，随机种子固定为 2026。

第 \(\omega\) 个情景为

\begin{equation}
L_{\tau,t}^{(\omega)}
=
\max
\left\{
0,
\hat L_{\tau,t}^{(d)}
+
e^L_{r_\omega,t}
\right\},
\end{equation}

\begin{equation}
PV_{\tau,t}^{(\omega)}
=
\max
\left\{
0,
\widehat{PV}_{\tau,t}^{(d)}
+
e^{PV}_{r_\omega,t}
\right\}.
\end{equation}

SAA 的作用是用有限个历史误差情景近似未来预测误差分布，使 0:00 的计划购电能够提前考虑 5 倍紧急购电风险。

\subsection{7 日滚动视野}

若仅优化当天，模型容易在当天结束前大量放电，从而忽略剩余储电量对未来日期的价值。为减小该末端效应，每天 \(d\) 的 0:00 同时优化

\[
d,d+1,\ldots,d+6
\]

共 7 天，但只真正执行第 1 天。

后 6 天的计划变量仅用于估计当前日末储电量的未来延续成本，并不代表今天提前提交未来日期的正式购电计划。

由此取滚动视野

\[
R=7.
\]

\subsection{随机情景下的物理约束}

对每个 \((\omega,\tau,t)\)，定义

\[
\overline L_{\tau,t}^{(\omega)}
=
\max
\left(
L_{\tau,t}^{(\omega)}
-
PV_{\tau,t}^{(\omega)},
0
\right),
\]

\[
\overline{PV}_{\tau,t}^{(\omega)}
=
\max
\left(
PV_{\tau,t}^{(\omega)}
-
L_{\tau,t}^{(\omega)},
0
\right).
\]

负荷平衡为

\begin{equation}
G^{L,(\omega)}_{\tau,t}
+
D^{(\omega)}_{\tau,t}
+
H^{(\omega)}_{\tau,t}
=
\overline L_{\tau,t}^{(\omega)}.
\end{equation}

充电来源为

\begin{equation}
PV^{ch,(\omega)}_{\tau,t}
+
G^{ch,(\omega)}_{\tau,t}
=
C^{(\omega)}_{\tau,t},
\end{equation}

且紧急购电 \(H\) 不用于储能充电。

光伏剩余平衡为

\begin{equation}
PV^{ch,(\omega)}_{\tau,t}
+
V^{(\omega)}_{\tau,t}
=
\overline{PV}_{\tau,t}^{(\omega)}.
\end{equation}

储能递推为

\begin{equation}
E^{(\omega)}_{\tau,t}
=
E^{(\omega)}_{\tau,t-1}
+
0.9C^{(\omega)}_{\tau,t}\Delta t
-
\frac{D^{(\omega)}_{\tau,t}\Delta t}{0.9}.
\end{equation}

每个滚动日开始时，所有情景共享当天真实储电量

\begin{equation}
E^{(\omega)}_{d,0}
=
E^{act}_{d,0}.
\end{equation}

容量、功率和互斥约束为

\[
1200\le E^{(\omega)}_{\tau,t}\le10800,
\]

\[
0\le C^{(\omega)}_{\tau,t},D^{(\omega)}_{\tau,t}\le5000,
\]

\[
C^{(\omega)}_{\tau,t}
\le5000u^{(\omega)}_{\tau,t},
\]

\[
D^{(\omega)}_{\tau,t}
\le5000(1-u^{(\omega)}_{\tau,t}),
\]

\begin{equation}
u^{(\omega)}_{\tau,t}\in\{0,1\}.
\end{equation}

\subsection{目标函数}

每天 \(d\) 的正式随机优化目标为

\begin{equation}
\min Z_d
=
C_d^{first}
+
\frac1K
\sum_{\omega=1}^{K}
Q_d^{(\omega)}.
\end{equation}

其中第一阶段全天计划购电费用为

\begin{equation}
C_d^{first}
=
\sum_t
\pi_tG^{plan}_{d,t}\Delta t.
\end{equation}

情景 \(\omega\) 下的补救及未来延续成本为

\begin{equation}
Q_d^{(\omega)}
=
\sum_t
5\pi_tH_{d,t}^{(\omega)}\Delta t
+
\sum_{\tau=d+1}^{d+6}\sum_t
\left[
\pi_tG_{\tau,t}^{plan,(\omega)}
+
5\pi_tH_{\tau,t}^{(\omega)}
\right]\Delta t.
\end{equation}

由此，模型在 0:00 同时权衡

\[
\boxed{
\text{提前正常购电成本}
\quad\longleftrightarrow\quad
\text{未来 5 倍紧急购电风险}
}
\]

并通过 7 日延续成本保留储能的跨日价值。

\subsection{实际执行规则}

0:00 求得当天

\[
G^{plan}_{d,1:144}
\]

后，计划在当天不再修改。

真实运行时依次执行：

\begin{enumerate}
\item 实际光伏优先满足负荷；
\item 当前计划购电补充负荷；
\item 仍不足时储能放电；
\item 储能仍不足时紧急购电；
\item 负荷满足后，富余光伏优先充电；
\item 再利用剩余计划电给储能充电；
\item 剩余光伏记为 \(V\)；
\item 剩余已购正常电记为 \(W\)；
\item 同一时段不得同时充放电。
\end{enumerate}

真实储电量按

\begin{equation}
E^{act}_{d,t}
=
E^{act}_{d,t-1}
+
0.9C^{act}_{d,t}\Delta t
-
\frac{D^{act}_{d,t}\Delta t}{0.9}
\end{equation}

更新，并将

\[
E^{act}_{d+1,0}
=
E^{act}_{d,T}
\]

传递至下一滚动日。

\subsection{模型求解}

经 SAA 情景展开后，模型的目标函数和能量平衡、储能递推、计划连接等约束均为线性形式，同时存在充放电状态二元变量 \(u\)，因此确定性等价模型属于混合整数线性规划。

采用分支定界/分支切割方法求解，并使用 MATLAB \texttt{intlinprog} 实现。首先放松整数约束求解 LP 松弛问题，并利用 LP 解进行热启动；随后对非整数状态变量进行分支，计算各节点下界，利用不可行剪枝和界剪枝排除不可能优于当前整数解的节点，同时使用切平面进一步收紧可行域。

每天的计算流程为

\[
\boxed{
\begin{array}{c}
\text{读取截至 }d-1\text{ 日历史数据}\\
\downarrow\\
\text{预测未来 7 日负荷与光伏}\\
\downarrow\\
\text{更新最近 28 个有效残差日}\\
\downarrow\\
\text{抽取 }K=4\text{ 个联合误差情景}\\
\downarrow\\
\text{构建 7 日滚动两阶段 SAA-MILP}\\
\downarrow\\
\text{LP 松弛热启动}\\
\downarrow\\
\text{分支定界/分支切割求解}\\
\downarrow\\
\text{提取当天 }G^{plan}\\
\downarrow\\
\text{按真实数据执行并更新 }E^{act}\\
\downarrow\\
\text{进入下一滚动日}
\end{array}
}
\]

全部算例正常收敛，最大约束违反约为

\[
4.07\times10^{-8},
\]

最大绝对最优间隙约为

\[
8.69\times10^{-4}\ {\rm 元},
\]

相对于全年总费用数值上可视为 0。

\subsection{求解结果与结果分析}

正式报送窗口为 2025 年 2 月 1 日至 12 月 31 日。

全天计划购电费用为

\[
13\,149\,064.98\ {\rm 元},
\]

紧急购电费用为

\[
1\,657\,825.30\ {\rm 元},
\]

总购电费用为

\[
\boxed{
14\,806\,890.28\ {\rm 元}
}.
\]

累计紧急购电量为

\[
269\,240.9\ {\rm kWh}.
\]

334 个报送日中，共有 175 天发生紧急购电，紧急购电时段共 1342 个，最大单个 10 min 时段紧急购电量为 966.770 kWh。

报送窗口日末储电量统计为

\[
\text{mean}=7046.0,
\qquad
\text{median}=8058.5,
\]

\[
\min=1200,
\qquad
\max=10800\ {\rm kWh}.
\]

四个指定日期结果如下：

\begin{table}[H]
\centering
\zihao{5}
\begin{tabular}{lrrrr}
\toprule
日期 & 日费用/元 & 紧急购电量/kWh & 日末储电量/kWh & 计划购电总量/kWh \\
\midrule
2025-03-20 & 41535.92 & 0.00 & 4077.0 & 66288.5 \\
2025-06-21 & 20827.44 & 0.00 & 8313.3 & 38911.0 \\
2025-09-23 & 71204.80 & 5351.69 & 5624.4 & 60974.9 \\
2025-12-21 & 66180.88 & 0.00 & 9391.2 & 101929.1 \\
\bottomrule
\end{tabular}
\caption{问题二典型日逐日指标（费用单位：元，电量单位：kWh）}
\label{tab:7}
\end{table}

%MANDATED:BEGIN
按题目表 1、表 2 与表 3 要求的格式，上述四个指定日期的结果整理如表~\ref{tab:m1-q2} 至表~\ref{tab:m3-q2}。

\begin{table}[H]
\centering
\zihao{5}
\begin{tabular}{lrrrr}
\toprule
指标 & 2025.3.20 & 2025.6.21 & 2025.9.23 & 2025.12.21 \\
\midrule
10:00-10:10 & 0.00 & 0.00 & 0.00 & 0.00 \\
12:00-12:10 & 730.72 & 0.00 & 330.01 & 970.04 \\
14:00-14:10 & 0.00 & 0.00 & 0.00 & 0.00 \\
16:00-16:10 & 507.64 & 68.68 & 424.60 & 847.50 \\
18:00-18:10 & 690.59 & 318.42 & 595.42 & 719.18 \\
20:00-20:10 & 0.00 & 0.00 & 0.00 & 0.00 \\
\midrule
全天购电量/kWh & 66288.55 & 38911.04 & 60974.88 & 101929.11 \\
全天购电费/元 & 41535.92 & 20827.44 & 71204.80 & 66180.88 \\
\bottomrule
\end{tabular}
\caption{微网在指定时间段的购电量及全天的购电量和购电费（问题二）}
\label{tab:m1-q2}
\end{table}
\par\vspace{1pt}{\zihao{6}注：购电量单位 kWh，购电费单位元；列对应题目表 3 的四个指定日期。}


\begin{table}[H]
\centering
\zihao{5}
\begin{tabular}{lrrrr}
\toprule
时间段 & 2025.3.20 & 2025.6.21 & 2025.9.23 & 2025.12.21 \\
\midrule
\multicolumn{5}{l}{\textit{充电量/kWh}} \\
0:00-4:00 & 1209.39 & 310.78 & 5158.75 & 2366.13 \\
4:00-8:00 & 1641.51 & 788.78 & 850.99 & 929.13 \\
8:00-12:00 & 5458.12 & 9897.19 & 2242.72 & 5251.08 \\
12:00-16:00 & 3881.82 & 0.00 & 6336.46 & 6329.48 \\
16:00-20:00 & 180.35 & 0.00 & 11.54 & 49.51 \\
20:00-24:00 & 3173.35 & 7987.49 & 4918.46 & 8405.67 \\
\midrule
\multicolumn{5}{l}{\textit{放电量/kWh}} \\
0:00-4:00 & 346.66 & 348.09 & 805.22 & 100.62 \\
4:00-8:00 & 4545.26 & 1724.85 & 6502.94 & 633.41 \\
8:00-12:00 & 3265.56 & 0.00 & 1140.14 & 7197.13 \\
12:00-16:00 & 259.22 & 0.00 & 682.05 & 2183.12 \\
16:00-20:00 & 5075.92 & 5584.97 & 6273.05 & 4461.69 \\
20:00-24:00 & 3686.76 & 3122.97 & 2.00 & 3654.96 \\
\midrule
0:00 储电量/kWh & 9175.12 & 3206.20 & 5174.47 & 8649.86 \\
24:00 储电量/kWh & 4077.01 & 8313.25 & 5624.39 & 9391.17 \\
\bottomrule
\end{tabular}
\caption{储能设备在指定时间段的充放电量及 0:00 和 24:00 的储电量（问题二）}
\label{tab:m2-q2}
\end{table}
\par\vspace{1pt}{\zihao{6}注：充放电量与储电量单位均为 kWh。}


\input{tables/tab3_q2}

%MANDATED:END


为分析预测不确定性的成本，构造“全年联合完美信息—逐日完美信息—7 日滚动 SAA”三级信息集阶梯。三种模型对应费用分别为

\[
12\,227\,142.86\ {\rm 元},
\]

\[
12\,253\,936.65\ {\rm 元},
\]

\[
14\,806\,890.28\ {\rm 元}.
\]

因此正式随机模型相对于全年联合完美信息增加

\[
2\,579\,747.42\ {\rm 元}
=
257.97\ {\rm 万元}.
\]

其中，跨日/有限视野部分为

\[
12\,253\,936.65
-
12\,227\,142.86
=
26\,793.79\ {\rm 元}
=
2.68\ {\rm 万元},
\]

预测不确定性部分为

\[
14\,806\,890.28
-
12\,253\,936.65
=
2\,552\,953.63\ {\rm 元}
=
255.30\ {\rm 万元}.
\]

即

\[
\boxed{
257.97
=
2.68+255.30
\quad(\text{万元})
}
\]

说明问题二中成本上升的主要来源是负荷和光伏预测不确定性，而不是滚动视野不足。

\subsubsection{滚动视野分析}

当

\[
R=1,3,7
\]

时，总费用分别为

\[
14\,884\,621.29,\quad
14\,806\,918.29,\quad
14\,806\,890.28\ {\rm 元},
\]

日末平均储电量分别为

\[
1961.4,\quad
7045.7,\quad
7046.0\ {\rm kWh}.
\]

由 \(R=1\) 扩大至多日视野后，日末储电量明显恢复，说明多日滚动能够有效体现储能跨日价值；而 \(R=3\) 到 \(R=7\) 的边际费用改善已经很小。

\subsubsection{情景数稳定性}

当 \(K=4\) 时，

\[
Z=14\,806\,890.28\ {\rm 元},
\qquad
Q_{em}=269\,240.9\ {\rm kWh}.
\]

当 \(K=8\) 时，

\[
Z=14\,516\,318.96\ {\rm 元},
\qquad
Q_{em}=230\,031.1\ {\rm kWh}.
\]

总费用差异约 1.96\%，但紧急购电量差异约 14.56\%，说明费用对情景数相对稳定，而紧急购电风险对情景覆盖仍较敏感。

\subsubsection{偏差校正分析}

不带偏差校正时，

\[
14\,671\,204.28\ {\rm 元},
\qquad
274\,289.2\ {\rm kWh}.
\]

带偏差校正时，

\[
14\,806\,890.28\ {\rm 元},
\qquad
269\,240.9\ {\rm kWh}.
\]

即总费用增加约 0.92\%，紧急购电量下降约 1.8\%。因此偏差校正不是单纯降低成本的模块，而是在经济性与供电可靠性之间进行折中。
